\section{Related Work}
\label{sec:related}

\paragraph{Membership inference and activation-level evidence.}
Classical membership inference attacks use shadow models~\citep{shokri2017membership} or loss thresholds~\citep{yeom2018privacy}, while LLM attacks use token probabilities, likelihood ratios, calibrated scores, and neighborhood comparisons~\citep{carlini2021extracting,shi2023detecting,zhang2025minkpp,ye2022enhanced,mattern2023membership,mireshghallah2022quantifying}. Fine-tuned-LLM attacks further refine these signals with local windows and in-context probing~\citep{chen2026wbc,lu2026icp}. Prior work also shows that benchmark artifacts such as temporal shift, vocabulary imbalance, candidate-set construction, and paraphrase sensitivity can dominate reported performance~\citep{duan2024membership,meeus2025sok}. A complementary white-box line finds membership evidence in internal states, including layer-wise hidden activations, attention and hidden-state statistics, gradient-induced activation drift, and pre-versus-post fine-tuning activation differences~\citep{galli2024lumia,ferrando2026memtrace,huang2026gdrift,minder2025narrow}. These works establish that activations can carry membership evidence. We ask how that evidence behaves under an SAE encode--decode decomposition.

\paragraph{Sparse autoencoders and incomplete reconstruction.}
Sparse autoencoders learn overcomplete dictionaries that express activations through sparse feature coordinates~\citep{cunningham2024sae,bricken2023monosemanticity}, building on the superposition hypothesis of \citet{elhage2022toy}. Recent work scales SAE training with TopK-style objectives and frontier-model dictionaries~\citep{gao2025scaling,templeton2024scaling}. At the same time, SAE work cautions against treating the learned dictionary as a complete account of the activation. SAE reconstruction error can contain structured components with downstream effects~\citep{engels2024dark}, and broader evaluations identify limits in reconstruction fidelity, feature completeness, feature canonicity, and steering-vector decomposition~\citep{karvonen2025saebench,leask2025sparse,mayne2024steering,sharkey2025open,peng2025discover}. Recent probing studies similarly find that SAE reconstructions can underperform residual stream or tuned linear baselines~\citep{smith2025deepmind,kantamneni2025sparse,davis2026probes}. We bring this concern to membership privacy by testing whether fine-tuning membership evidence is preserved by the SAE reconstruction or remains recoverable from the reconstruction residual.

\paragraph{SAE-based privacy interventions and unlearning.}
A growing line of work uses SAE features as the substrate for targeted unlearning, guardrails, and privacy interventions. SSPU projects weights onto SAE-defined subspaces~\citep{wang2025sspu}. PrivacyScalpel ablates SAE features to suppress personally identifiable information leakage~\citep{frikha2025privacyscalpel,grande2026privacyscalpelv3}. DSG uses dynamic SAE guardrails for inference-time intervention~\citep{muhamed2025dsg}. \citet{suri2025mitigating} apply SAE feature steering to memorization mitigation in Gemma, and \citet{farrell2024applying} report a related negative result on WMDP-Bio unlearning. These methods act through SAE feature coordinates, while other memorization-mitigation approaches act on raw activations or circuit-level primitives~\citep{zhang2026gss,lasy2025memcircuits}. Our audit tests whether membership evidence is actually preserved by the representation that feature-level interventions edit. If detectors recover membership signal from the reconstruction residual, then suppressing SAE features or extractable outputs does not by itself demonstrate reduced membership inference risk.

\paragraph{Contrastive geometry and robustness controls.}
Our analysis also connects to work that uses feature patching and contrastive geometry to localize model functions. Sparse feature circuits localize functions by patching SAE features~\citep{marks2024sparse}, while contrastive principal component analysis, representation engineering, and activation addition identify directions associated with model outputs~\citep{abid2018exploring,zou2023representation,rimsky2024steering}. Some functions can be mediated by a single residual stream direction~\citep{arditi2024refusal}, but our residual-PC intervention is more limited. It tests partial concentration of membership signal and is not treated as complete causal localization. Following the logic of double dissociation from neuropsychology, we treat membership detection and verbatim extraction as potentially separable outcomes, where one may be present without the other, implying that they should not be interpreted as the same privacy signal~\citep{teuber1955physiological,shallice1988neuropsychology}, and include perturbation and low-FPR controls motivated by known fragilities in membership inference evaluation~\citep{duan2024membership,meeus2025sok}.


